\section{An explicit cyclotomic depth for the two quartic hosts}
\label{sec:cubic-effective-host-depth}

This proof replaces the ineffective high-layer choices in the actual
quartic host constructions.  It keeps the original weights, the full
local groups, and every retained automorphic and native-source input.
The finite class certificates below are quotient certificates; no
complete unit basis or full class-number equality is asserted.

\paragraph{Fixed arithmetic and notation.}
Put $p=269$, $K=\mathbb Q(\sqrt{-3})$, and $\eta=\chi_{-3}$.
Let $\Psi$ be the quartic character of conductor five with
$\Psi(2)=i$, where $i=82$ in $\mathbb F_p$.
The symbol $c_p$ denotes the actual action on $\mu_p$, and $\omega$
its Teichmuller lift.  Work over any finite field $k$ containing all
the residual characters, and put
\[
 \gamma=\Psi c_p^e,\quad e\equiv267\text{ or }133\pmod{268},
 \qquad \alpha=\chi^{-j}\gamma^{-1},\quad j=1,2.
 \tag{G140.1}
\]
The character $\alpha$ is on the CM field; the other displayed
characters extend from $\mathbb Q$.  Retain the actual first- or
second-residue construction of AC-G136 or AC-G138.  Its initial fields
and two cyclic steps are
\[
 \begin{gathered}
 F_{\rm ab}=\mathbb Q(\sqrt{33},\sqrt2),\qquad
 H_{\rm ab}=KF_{\rm ab},\qquad [F_{\rm ab}:\mathbb Q]=4,\\
 F_1=F_{\rm ab}T_\theta,\quad H_1=H_{\rm ab}T_\theta,
 \qquad F_0=H_0^+,\quad H_0=H_1T_q,\\
 [F_1:F_{\rm ab}]=[F_0:F_1]=p,
 \qquad [F_0:\mathbb Q]=4p^2.
 \end{gathered} \tag{G140.2}
\]
Here $T_\theta/\mathbb Q$ is the actual real cyclic degree-$p$
field of conductor $q_\theta=16141$, so $v_p(q_\theta-1)=1$.
The actual AC-G83 field $T_q/\mathbb Q$ is dihedral with quadratic
field $E=\mathbb Q(\sqrt2)$ and $[T_q:E]=p$.
Its prime $q$ is split in $K\mathbb Q(\sqrt{-11})T_\theta$,
inert in $E$, and satisfies $v_p(q+1)=1$.
Thus it has residue degree two in $F_1$ and in $H_1$.
All these fields are unramified at $p$.  The completions of
$F_{\rm ab},H_{\rm ab}$ have degree two, and those of
$F_1,H_1,F_0,H_0$ have degree $2p$, over $\mathbb Q_p$.

We refine the AC-G92 prime choice as proved below, using one selected
real place.  Let $T'/F_0$ be its cyclic extension, with
\[
 [T':F_0]=p^s\le p h_{F_0,S}(p),\quad s\ge1,\qquad
 F_{\rm out}=T',\quad H_{\rm out}=H_0T',\quad m=p^{s+2}.
 \tag{G140.3}
\]
The notation $T'$ here includes $F_0$, as in AC-G92.
For $n\ge0$, append the degree-$p^n$ cyclotomic layer and write
$F_n,H_n$.  Their real degree and local degrees are
\[
 d_n=4mp^n,\qquad d_{v,n}=2p^{n+1}.
 \tag{G140.4}
\]
All Iwasawa modules use the same coordinate $T=\sigma-1$, where
$\sigma$ is a generator of $\operatorname{Gal}(\mathbb Q_\infty/
\mathbb Q)$.  In particular $T^{p^n}$ is the characteristic-$p$
coinvariant ideal at the actual layer $n$; no reindexing is made.

\begin{theorem}[AC-G140: effective depth for the actual quartic construction]
\label{AC-G140}
The selected prime in (G140.3) can be chosen with
$v_p(q_*-1)=2$, while retaining every AC-G92--AC-G93 local character,
class-positivity and host property.  On this host, put
\[
 C_0=m(12+2/p),\qquad A=p^s(p^2+2p),\qquad n_*=s+4.
 \tag{G140.5}
\]
These are explicit integers.  Write $a_n$ for the strict adjacent
dimension, $z_n$ for the strict central dimension, and $h_n$ for the
full trivial degree-one dimension in AC-G98.  Then, for every $n\ge0$,
\[
 \begin{split}
 a_n&\le m(3p^n+12+2/p),\\
 z_n&\le m(6+1/p),\qquad h_n\le1+m(6+1/p).
 \end{split} \tag{G140.6}
\]
The literal rational-frame dimensions satisfy
$\mathrm{b1}\le C_0$, $\mathrm{b1p}\le C_0$, and
$\mathrm{h1}\le1+C_0$.  For every $n\ge1$ both
\[
 h_n+a_n+z_n+1<d_n,\qquad
 3C_0+7<d_n,\quad d_{v,n}>5
 \tag{G140.7}
\]
hold.  The rational compact modules are annihilated by $T^A$.
Restriction on the strict mixed obstruction is zero from layer
$s+3$ to layer $s+4$.  Thus $n_*$ is a common explicit layer for
the nonzero filled Low classes, the full rational
$\mathrm{Frame.Ok}/\mathrm{Frame.Big}$, and the changed-tuple
dimension inequality used in AC-G136 and AC-G138.

All later native support, component, original-model and counting
arguments of those theorems may use this particular layer, with
the same retained hypotheses and the original coefficient DVR
unchanged.  The source seed weight stays $e_*=267$ or $133$,
respectively.  This is an effective depth relative to the actual
selected extension; it supplies no global character-box budget or
mixed theta/Selmer realization.
The same bound applies to each fresh signed and compensated host
construction in AC-G139, with its own selected extension and native
ideals, as verified below.
\end{theorem}

\begin{proof}
\emph{1. Exact finite transfer before Iwasawa localization.}
Use the setup of AC-G100: $L/B$ has prime-to-$p$ group $\Delta$,
$N=LT$ with $T/B$ cyclic of order $b=p^u$, and the finite extension
is disjoint from $B_\infty/B$.  Let $X'_J$ denote the maximal
unramified, $p$-split abelian pro-$p$ module over $J_\infty$.
For a nontrivial one-dimensional $\vartheta$ on $\Delta$, set
$M_J=e_\vartheta((X'_J/p)\otimes k)$.
The compact commutator calculation of (G100.10) gives
\[
 0\longrightarrow (X'_N)_P\longrightarrow\mathscr G^{\rm ab}
       \longrightarrow P\longrightarrow0.
 \tag{G140.8}
\]
Both $P[p]$ and $P/p$ in its reduction modulo $p$ have trivial
$\Delta$ action.  Their $\vartheta$-parts vanish.  Consequently
\[
 (M_N)_P\simeq e_\vartheta((\mathscr G^{\rm ab}/p)\otimes k)
 \tag{G140.9}
\]
before any localization.  Killing the finite subgroup $R$ generated
by away-$p$ inertia and the full $p$-decomposition images gives
exactly $X'_L$, by (G100.13).  Hence the kernel of the map from
(G140.9) onto $M_L$ is an image of $e_\vartheta((R/p)\otimes k)$.
After quotienting by $T^{p^n}$, right exactness and the nilpotent
augmentation ideal of $k[P]$ imply
\[
 \dim_k M_N/T^{p^n}M_N
 \le b\left(\dim_k M_L/T^{p^n}M_L+C_\vartheta\right),
 \quad C_\vartheta=\dim_k e_\vartheta((R/p)\otimes k).
 \tag{G140.10}
\]
Indeed a basis of coinvariants lifts to generators: if its span
generates a submodule $U$, then $V/U=J(V/U)$, so $J^b=0$ gives
$V/U=0$.  Thus $\dim V\le b\dim V_P$.

We also retain the operator bound.  If $T^aM_L=0$ and the error
is killed by $T^c$, then $T^{a+c}(M_N)_P=0$, whence
$T^{a+c}M_N\subset JM_N$.  Iterating $b$ times gives
\[
                         T^{b(a+c)}M_N=0.
 \tag{G140.11}
\]
This argument requires neither semisimplicity of $P$ nor a
vanishing Iwasawa invariant.

\emph{2. The precise ramification errors and the initial control.}
Every local image generating $R$ injects into the cyclic group $P$.
Its stabilizer acts trivially: after injection that conjugation is
conjugation on the commuting cyclic quotient.  Thus $R/p$ is a
quotient of a sum of prime-permutation modules, one generator per
relevant prime.  The $p$-part of this permutation sum has zero
$\vartheta$-part whenever $\vartheta$ is nontrivial on every
$p$-decomposition group.  For each away-$p$ prime $v$ of $B$, a
$\Delta$-orbit contributes at most one dimension to a given
one-dimensional character, for each cyclotomic orbit.
If $c_v$ is the number of primes above $v$ in $B_\infty$, then
\[
 C_\vartheta\le\sum_{v\text{ ramified away }p}c_v,
 \qquad T^{\max_v c_v}\,e_\vartheta(R/p\otimes k)=0.
 \tag{G140.12}
\]
Here $c_v$ is a power of $p$ and its permutation module is
$k[T]/T^{c_v}$.  Since our cyclotomic coordinate is the rational
one, a residue norm $Nv$ gives
\[
       c_v=p^{v_p((Nv)^{p-1}-1)-1}.
 \tag{G140.13}
\]
To see this, the principal-unit part of $Nv$ has logarithm of
valuation $v_p((Nv)^{p-1}-1)$; a generator $1+p$ has logarithm
of valuation one.  Their index in $\mathbb Z_p$ is (G140.13).
There is no hidden character twist in this permutation action.
In the commuting product $\Delta\times\Gamma$, the decomposition
subgroup contains the pure open $\Gamma$ subgroup of that index:
its possible graph map into the finite prime-to-$p$ quotient of
$\Delta$ is a homomorphism from a pro-$p$ group and is trivial.
Projecting to a one-dimensional $\Delta$ character therefore
either removes that summand or leaves the stated $\Gamma$ module.

The common residual splitting covers at layer zero are tamely
ramified at $p$.  They are composites of prime-to-$p$ coefficient
extensions with the unramified $p$-power completions in (G140.2).
The selected extension splits at $p$.  A finite tame extension
of $\mathbb Q_p$ is disjoint from its totally ramified cyclotomic
$\mathbb Z_p$-extension, and all the latter steps stay totally
ramified after that base change.  Therefore the control sequence
AC-G74 applies from layer zero, without reindexing.

It is the actual coefficient $\vartheta$, not its later Kummer
twist, which enters that control sequence.  For $\alpha$, its
cubic tame factor of order three cannot cancel the rational
factor of order dividing $268$.  For $\gamma^{-1}$ and
$\eta\gamma^{-1}$, the inertia character is $c_p^{-e}$ of
order $268$; for $W=\eta\gamma^{-2}$ it is $c_p^2$ of order
$134$; and for $c_p$ it has order $268$.
All these restrictions remain nontrivial at every $p$-place
after all the $p$-power steps.  Thus each corresponding
$p$-permutation correction in AC-G74 vanishes, and
\[
 \dim H^1_{\rm str}(B_n,k\vartheta)
                 =\dim_k M_B/T^{p^n}M_B.
 \tag{G140.14}
\]
These are also the literal native away-kernel groups.  To compare
with a smaller coefficient splitting field, restrict a native
cocycle there, where the action is trivial.  Its additive restriction
kills an away-inertia subgroup of prime-to-$p$ index, hence the whole
away-inertia group.  Conversely the unramified class descends by
exact prime-to-$p$ invariants.  At $p$ use every place and the full
local groups; restriction and corestriction with prime-to-$p$ index
give the same strict zero condition.  This is the exact comparison
of AC-G90/AC-G93, not an assumed inflation injection on $H^2$.
The same assertion holds for the CM adjacent group.
In particular the existing actual certificate $a_{\rm ab}(0)\le3$
gives, by $k[[T]]$-Nakayama,
\[
                         a_{\rm ab}(n)\le3p^n.
 \tag{G140.15}
\]
No assertion about the length of an unspecified cubic torsion
module is needed.

\emph{3. Initial rational strict groups vanish.}
We prove, with their actual full local conditions,
\[
 \begin{gathered}
 H^1_{\rm str}(H_{\rm ab},k\gamma^{-1})=0,\\
 H^1_{\rm str}(F_{\rm ab},kW)=
 H^1_{\rm str}(F_{\rm ab},kc_p)=0,
 \qquad W=\eta\gamma^{-2}=\chi_{-15}c_p^2.
 \end{gathered} \tag{G140.16}
\]
The first line is, by quadratic Shapiro, the sum of the two real
strict groups with coefficients $\gamma^{-1}$ and
$\eta\gamma^{-1}$.  Their full $p$-local summands are included.

For that first line, Shapiro from $H_{\rm ab}/K$ gives the four
quadratic twists $1,\tau_{-11},\tau_8,\tau_{-88}$.
For a coefficient $\beta=\gamma^{-1}\delta$ its Kummer character
is $\beta/c_p=\Psi^{-1}c_p^{-(e+1)}\delta$.
It is cut out over $K$ by $M_d=KF_d$, where
\[
 F_d=\mathbb Q(w),\quad w^4+5dw^2+5d^2=0,\qquad
 d\in\begin{cases}
 \{1,-11,2,-22\},&e\equiv267,\\
 \{269,-2959,538,-5918\},&e\equiv133.
 \end{cases} \tag{G140.17}
\]
The inverse Kummer eigenspace is primitive of order four;
both primitive characters will be checked.  Since $F_d$ has
unique quadratic subfield $\mathbb Q(\sqrt5)$, the group of
$M_d/\mathbb Q$ is $C_4\times C_2$.
Its four characters primitive on $\operatorname{Gal}(M_d/K)$
come from the quartic fields $F_d$ and $F_{-3d}$.
The exact flag-one certificates give the following class-quotient
orders, all prime to $p$:
\[
\begin{array}{c|rrrrrrrr}
d&1&-11&2&-22&269&-2959&538&-5918\\\hline
F_d&1&1&2&2&68&4&272&16\\
F_{-3d}&1&8&2&64&8&1600&40&21760
\end{array} \tag{G140.18}
\]
For each character, norm to its quartic field followed by extension
back to $M_d$ is multiplication by two on that character
sector.  Therefore all primitive order-four class sectors in
$M_d$ vanish.  This does not assert that every other class
sector of $M_d$ vanishes.

A strict Kummer class is unramified at every finite prime away
from $p$, and zero at every prime above $p$.  Its valuations
at all finite primes are consequently multiples of $p$.
The Kummer ideal sequence and (G140.18) reduce its primitive
sector to ordinary global units, with no $S$-divisor error.
The CM unit quotient by the units of the maximal real subfield
is killed by two after removing roots of unity: $u/\bar u$ has
all conjugates of modulus one and is a root of unity, while
$u^2=(u\bar u)(u/\bar u)$.
Moreover
there is no $p$-torsion root of unity in these degree-eight fields.
The relevant real cyclic quartics have parameters
\[
             r=-3,-11,-6,-22,-807,-2959,-1614,-5918.
 \tag{G140.19}
\]
Their two primitive unit eigenspaces each have dimension one,
by the equivariant Dirichlet unit theorem for a real cyclic
quartic field.  The following supplied units, in the basis
$1,w,w^2,w^3$ for $F_r$, have monic integral characteristic
polynomials with constant term $\pm1$:
\[
\begin{array}{c|rrrr}
r&a_0&a_1&a_2&a_3\\\hline
-3&1/2&1/2&-1/6&0\\
-11&-4&3/2&1/22&-1/22\\
-6&-7&2&2/3&-1/6\\
-22&51&8&-8/11&-1/11\\
-807&971/2&-29/2&-1/6&4/807\\
-2959&-2479374&77545/2&5095/22&-21433/5918\\
-1614&3953&-84&-2/3&23/1614
\end{array} \tag{G140.20}
\]
For the last field the supplied unit is
\[
 \begin{aligned}
 u_{-5918}={}&2001597290961-22133013208w\\
 &+\frac{-276613942432w^2+3058707198w^3}{2959}.
 \end{aligned}
\]
All these finite assertions are checked by multiplication
matrices in $\mathbb Q[w]/(w^4+5rw^2+5r^2)$ in the exact
verifier \path{degree4/verify_units.py}; it prints every monic
relation, not just numerical approximations or norms.

The order-four automorphism is $b(w)=-(3+w^2/r)w$.
For the first four fields, the algebra at $p$ is unramified
with residue degrees at most two.  Compute
$\ell=(u^{p^2-1}-1)/p$ in its integral algebra modulo $p$.
The two unnormalized projectors $\sum_{a=0}^3(\pm i)^{-a}b^a\ell$
have coefficients $(0,A,0,B)$ as follows:
\[
\begin{array}{c|rr|rr}
r&A_+&B_+&A_-&B_-\\\hline
-3&77&6&182&13\\-11&216&151&118&163\\
-6&24&235&96&234\\-22&225&17&129&72
\end{array} \tag{G140.21}
\]
For the remaining fields put $R=w^2/r$; its residues are $69,195$
and $R^2+5R+5=0$.  Now $w$ is a ramified quadratic uniformizer.
Rewrite $u=A(R)+wB(R)$ with $p$-integral coefficients before
reducing.  The linear coefficient of $u^{p-1}$ is $-B/A$.
Projecting it with $b(R)=-5-R$, $b(w)=-(3+R)w$ gives
\[
\begin{array}{c|rr|rr}
r&(\ell_+)_{69}&(\ell_+)_{195}&(\ell_-)_{69}&(\ell_-)_{195}\\\hline
-807&54&73&157&261\\-2959&25&248&220&131\\
-1614&240&175&164&127\\-5918&258&20&75&63
\end{array} \tag{G140.22}
\]
Each vector is nonzero.  A local $p$-th power has zero indicated
logarithm, or zero uniformizer-linear coefficient in the ramified
case.  Thus localization is nonzero, hence injective, on both
one-dimensional primitive unit sectors.  The further local
extension to $M_d$ has degree at most two; restriction of
Kummer classes remains injective by restriction--corestriction.
This proves the first line of (G140.16), with both character
orientations and all local places accounted for.

For the central coefficient the Kummer twist is
$W/c_p=\chi_{-15}c_p$, of potentially large order.  We do not
replace it by a quartic field.  The four Shapiro characters over
$\mathbb Q$ are $\chi_D c_p^2$, with
$D=-15,-55,-120,-440$.  In the primitive cyclotomic field
$E_D=\mathbb Q(\mu_{|D|p})$, Stickelberger's theorem says that
$(\sigma_b-b)\Theta_{E_D}$ annihilates the class group, where
\[
 \Theta_{E_D}=\frac1{|D|p}
       \sum_{(a,|D|p)=1}a\sigma_a^{-1}.
 \tag{G140.23}
\]
Its scalar on the $\chi_D\omega^2$ sector is
$(\chi_D(b)\omega(b)^2-b)B_{1,\chi_D\omega^{-2}}$.
Exact sums modulo $p^2$, divided by $p$ only after their numerator
is verified divisible by $p$, give
\[
\begin{array}{c|rrrr}
D&-15&-55&-120&-440\\\hline
B_1\bmod p&248&1&260&249\\
b&2&2&7&3\\
(\chi_D(b)b^2-b)B_1\bmod p&227&2&235&149
\end{array} \tag{G140.24}
\]
Every annihilator scalar is a unit.  The relevant class sectors
therefore vanish.  Passage to a common cyclotomic field uses
extension and norm, not an imprimitive Stickelberger value:
on the inflated character sector their composite is the
prime-to-$p$ extension degree.  For example
$\mathbb Q(\mu_{1320p})$ contains all the needed fields and
$F_{\rm ab}$.  This avoids any omitted or zero Euler factor.
Explicitly, on the larger-field sector it is extension after norm,
$\iota\circ N=\sum_g g=[E:E_D]$, that proves vanishing.
Exact finite-prime-to-$p$ descent and the class convention
of AC-G74 prove the central strict vanishing in (G140.16).

For $c_p$ the four real Shapiro factors are $\delta c_p$ with
$\delta=1,\chi_{33},\chi_8,\chi_{264}$.  The same primitive
calculation gives, respectively, annihilator residues
\[
                  1,224,161,6\quad\text{in }\mathbb F_{269}.
 \tag{G140.25}
\]
For $\delta\ne1$, use $B_{1,\delta\omega^{-1}}
=139,18,38$ modulo $p$ and $b=5,3,7$.
For $\delta=1$ use $b=1+p$:
$(\sigma_{1+p}-(1+p))\Theta=-p\Theta$, whose scalar is
$-\sum_{a=1}^{p-1}a\omega(a)^{-1}=1\pmod p$.
This handles the one possibly nonintegral Bernoulli value
without dividing by $p$ incorrectly.
It proves the final strict zero in (G140.16).

By (G140.14), each zero in (G140.16) says $M/TM=0$ for its
actual coefficient module.  Compact Nakayama gives $M=0$.
Consequently all four real rational strict modules are zero
on the entire initial abelian cyclotomic tower.

\emph{4. A selected prime with exactly two powers of $p$.}
In AC-G92 choose the common finite Galois cover $M/F_0$ to
contain $\mu_p$ and all residual coefficients.  It remains
tamely ramified at $p$, so
$M\cap\mathbb Q(\mu_{p^\infty})=\mathbb Q(\mu_p)$.
Let $\mathcal K=M(E_{M,S}^{1/p})$.  This is finite, and
$\operatorname{Gal}(\mathcal K/M)$ is elementary abelian.
It is Galois over $\mathbb Q$: the full $S$-unit group is
stable under the Galois group, and roots of a generating set
already include roots of every conjugate unit.
Since $\mu_p\subset E_{M,S}$, one has $\mu_{p^2}\subset\mathcal K$.
But $M(\mu_{p^3})/M$ is cyclic of degree $p^2$, so it cannot
be contained in the exponent-$p$ extension $\mathcal K/M$.
It follows that
\[
 \mathcal K\cap\mathbb Q(\mu_{p^3})=\mathbb Q(\mu_{p^2}).
 \tag{G140.26}
\]
Identity on $\mathcal K$ and $1+cp^2$ on $\mu_{p^3}$, with
$c\not\equiv0\pmod p$, are therefore compatible automorphisms
of their compositum.  Chebotarev gives a fresh rational prime
$q_*$ in that class.  It splits in $M$, every $S$-unit is a
local $p$-th power there, and $v_p(q_*-1)=2$ exactly.
Choose one real place above it.  The unchanged AC-G92 ray
character then has inertia exactly $p$, is trivial at every
$p$-place, and has the degree bound (G140.3).
No faithful character of the bottom order-$p$ subgroup of a
residue group is assumed: the character is the quotient
character of order $p$, exactly as in AC-G92.  Thus the
AC-G93 inertia quotient by $pI$, positivity and local-norm
arguments apply without alteration.

\emph{5. Explicit dimensions and annihilators on the output host.}
For a real rational sector, the first extension has at most four
ramified base primes and all $c_v=1$, hence $C_1\le4$.
For the second, $q$ has residue degree two and splits in
$T_\theta$, giving $2p$ real base primes.  Since
$v_p(q^2-1)=1$, all have $c_v=1$, so $C_2\le2p$.
For the third step there is one selected real prime, of norm
$q_*$, and (G140.13) gives $c_v=p$, so $C_3\le p$.
Starting from the rational zero modules in Step 3, (G140.10) gives
\[
 \dim M_1/T^{p^n}\le4p,\quad
 \dim M_2/T^{p^n}\le6p^2,\quad
 \dim M_{\rm out}/T^{p^n}\le m(6+1/p).
 \tag{G140.27}
\]
The identical calculation over the CM base doubles the three
prime counts: $8,4p,2p$.  Starting with (G140.15) gives the
adjacent bound in (G140.6).

For the rational operator bound, the three error modules are
killed by $T,T,T^p$, respectively.  Applying (G140.11) gives
\[
       T^pM_1=0,\quad T^{p^2+p}M_2=0,\quad T^AM_{\rm out}=0.
 \tag{G140.28}
\]
These claims concern the actual nontrivial coefficient sectors;
the trivial $P[p]$ errors in (G140.8) have already been removed.

The source $\mathrm{b1}$ is the sum of the real
$\gamma^{-1}$ and $\eta\gamma^{-1}$ strict sectors by
quadratic Shapiro.  Hence it is at most $C_0$.
The central $\mathrm{b1p}$ has the smaller bound (G140.27).
For the full trivial group, global Euler characteristic and
compact-support duality give
\[
 h_n=1+\dim H^1_{\rm str}(F_n,kc_p).
 \tag{G140.29}
\]
The local invariant spaces of $kc_p$ vanish, since the actual
local fields do not contain $\mu_p$; so there is no local
degree-zero correction.  This proves the remaining bounds.
The estimates imply (G140.7) already for $n\ge1$: after division
by $m$, the two required strict inequalities follow from
\[
 p^n>24+4/p+2/m,
 \qquad 4p^n>36+6/p+7/m.
 \tag{G140.30}
\]
They hold for $p=269$, $n\ge1$, $s\ge1$.

\emph{6. A quantitative mixed-cup filling step.}
The exact dual character is
$\lambda=c_pW^{-1}=\chi_{-15}c_p^{-1}$.
It is even, nontrivial globally and on every $p$-inertia group.
Its global Euler characteristic is zero.  Since global
cohomological dimension is two and $H^0(\lambda)=0$, this gives
$\dim H^1(\lambda)=\dim H^2(\lambda)$.
The full arithmetic compact-support duality of AC-G90 and
AC-G96 identifies $H^2(\lambda)^\vee$ with $H^1_c(W)$.
All local $H^0(W)$ vanish, so the latter is the strict group.
Therefore at every layer
\[
                  \dim H^1(\Gamma_n,k\lambda)=z_n.
 \tag{G140.31}
\]
By (G140.28) the right side is constant once $p^n\ge A$.
Since $p^{s+3}>p^s(p^2+2p)$, this is true for every
$n\ge s+3$.  Restriction on the full $\lambda$ group is
injective by inflation--restriction: its nontrivial prime-to-$p$
image persists on an index-$p$ subgroup.  It is thus an
isomorphism at these layers.  Corestriction composed with
restriction is multiplication by $p$, hence zero over $k$.
Corestriction on the entire next-layer group is consequently zero.
The exact restriction/corestriction duality of AC-G90--AC-G91
now makes restriction on the central anti strict $H^2$ zero.

AC-G93 gives the required nonzero adjacent and central classes
at $H_{\rm out}$, as well as the rational positive classes
$\eta\gamma^{-1}$ and $\eta\gamma^{-2}$.
Their global and local invariant spaces vanish.  Restriction
through the $p$-power tower preserves their nonzero strict
classes.  The mixed obstruction is natural under restriction,
so at layer $s+4$ it vanishes and the actual polarized filling
and polynomial Low paths of AC-G90 exist.

\emph{7. The literal host and source interfaces.}
All nonnumerical $\mathrm{Frame.Ok}$ clauses were checked for
these actual fields in AC-G136/AC-G138: the full away-inertia
quotient, full local maps, parity and scalar witnesses,
$\mathrm{Psi\_N}$, the arbitrary-module Euler statement,
$\mathrm{PhiCond}$, and the local generators and global class
exponent of AC-G99.  They persist through the cyclotomic
$p$-extensions for the precise reasons given there.
The two positive rational classes above give $\mathrm{f4}$
and $\mathrm{f4'}$ at this particular layer.  Step 5 supplies
all three upper bounds with the same explicit $C_0$, and
(G140.7) is exactly the literal $\mathrm{Big}$ formula
$3C_0+5+2<d$, $5<d_v$.

The CM extension remains finite-unramified and relatively
split at all coefficient-prime and prescribed old-type places.
Both actual seeds remain irreducible and cuspidal through the
cyclic $p$-steps; ordered local flags and full finite torus
characters restrict exactly.  The auxiliary source prime is
chosen anew as before.  No condition of being unramified at
the coefficient prime is imposed by that source's distinct
auxiliary-prime hypothesis.  Thus the fresh native rings and
ideals in AC-G136/AC-G138 can be constructed at the fixed
layer $n_*$, retaining all original coefficient and selected
local-component conditions.
For AC-G139 replace $\Psi$ by $\Psi^{-1}$ if required.
The adjacent initial certificate is the same one, as proved in
its Step 3, and the two primitive quartic unit projectors were
both checked in (G140.21)--(G140.22).  The central coefficient and
the full trivial correction are unchanged because they use
$\Psi^2$.  The quadratic factor $\tau_8$ dies on the retained
host, and the factor $\tau_p^b$ is exactly compensated by
$r\equiv e+134b\pmod{268}$.  All actual coefficient modules
and their full local bounds are consequently those just proved,
or their already checked quartic inverses.  Apply the construction
branchwise to its own exact-root prime $q$ and selected extension.
The new native localization for the sign and each original
selected-component transport remain precisely those of AC-G139;
no identification of the different native ideals is required.
This proves the theorem at its stated retained-input scope.
\end{proof}

\paragraph{A bound not requiring the exact selected degree.}
The integer $s$ in (G140.3) records the $p$-power degree of the actual
ray-character image.  If only the old host is specified, put $N=4p^2$,
$P=2\cdot3\cdot11\cdot q_\theta q$, and
$B=(NP)^{N/2}$.  The local different estimate already proved
in AC-G86 gives $|D_{F_0}|\le(NP)^N$.
Minkowski and the coefficientwise bound $a_{F_0}(n)\le d_N(n)$
give
\[
 h_{F_0,S}(p)\le h_{F_0}
       \le B(1+\log B)^{N-1}.
 \tag{G140.32}
\]
Indeed each class has an ideal of norm at most $\sqrt{|D_{F_0}|}$,
and summing the $N$-fold divisor function is bounded by
$B(\sum_{r\le B}1/r)^{N-1}$.
Hence one may replace $s+4$ by the explicit integer
\[
 5+\left\lceil\log_p\bigl(B(1+\log B)^{N-1}\bigr)\right\rceil.
 \tag{G140.33}
\]
This bound depends on the already chosen actual prime $q$,
not on the sizes of the later selected primes.  Substituting
any previously retained effective bound for $q$ is legitimate;
no such substitution or global counting budget is asserted here.
In particular $s=\log_p[T':F_0]$ is part of the actual selected-host
datum, with $s\le1+v_p h_{F_0,S}$; the prime $q_*$ is still chosen
qualitatively.  No absolute numerical auxiliary-prime size or
field-discriminant bound is claimed.

\paragraph{Source and computation scope.}
The finite transfer and control arguments are new deductions
from the complete AC-G74/AC-G100 exact sequences; the proof
above retains all finite errors and checks the starting layer.
The published input for (G140.23) is Sharifi, \emph{Iwasawa
Theory}, Definition 4.5.1, Lemma 4.5.2 and Theorem 4.5.6.
The statements and the complete printed proofs of the lemma and
theorem were read in the author's chapter-four source.  Global reciprocity, Kummer theory,
Dirichlet's unit theorem, compact-support duality, and the
automorphic/native results retain the audited scopes in
AC-G74, AC-G90--AC-G100 and AC-G136--AC-G139.
The literal \path{ContinuousCoboundaries2.lean} fields
$\mathrm{Ok}$ and $\mathrm{Big}$ were reread; no Lean build
or new proof of the imported duality or automorphy theorems
is claimed.  Exact inputs, outputs, binaries, characteristic
polynomials, local projector verifiers and character-sum
scripts are preserved with the claim manifest.

\paragraph{Geometric dictionary and remaining scope.}
The ray steps are finite covers, etale away from their stated
ramification, with the selected local character and full
coefficient-prime splitting specified.  The cyclotomic step
changes that actual cover by its stated totally ramified local
degree.  The strict classes, their cup filling, and the native
ordinary points are those already constructed on these
arithmetic fundamental groups.  This uses the retained
finite-cover dictionary, not a new motive or an archival scan.
No simultaneous growing character-box estimate, all-prime
Selmer bound, or mixed theta/Selmer realization is supplied.
